% Work-MAT | 2. rocnik | pripravovana hodina 34
% Návaznost: uživatelské připravy.zip (hodiny 01–20), autorská učebnice na GitHubu.
% Realisticky.cz: didaktická inspirace v posloupnosti a typech úloh; zadání jsou vlastní.
\documentclass[11pt,a4paper]{article}
\input{../../sablony/pripravy/priprava.sty}
\nastavPripravu{2. ročník}{Logaritmy}{34}{Věty o logaritmech}
\begin{document}
\blok{Motivace a návaznost}
Začni rovností $8\cdot4=32$ a mocninami dvojky. Žáci mají objevit, že součinu argumentů odpovídá součet logaritmů, nikoliv logaritmus součtu exponentů.
\poznamka{Orientační nabídka pro 45--55 minut: krátké zopakování a motivace, nový výklad, tři společné vzorové úlohy, samostatná práce. Vyber podle tempa třídy; zbývající zadání slouží jako rezerva.}
\blok{Výklad}
Pro $u,v>0$ platí $\log_a(uv)=\log_a u+\log_a v$, $\log_a(\frac uv)=\log_a u-\log_a v$, $\log_a(u^r)=r\log_a u$.
\blok{Společné příklady}
\begin{prikladReseniBox}[title={Příklad 1}]
Vyřeš, resp. urči: $\log_2 8+\log_2 4$.
\tcblower
\[\begin{aligned}\log_2(8\cdot4)=\log_2 32 \\ =5\end{aligned}\]
\end{prikladReseniBox}
\begin{prikladReseniBox}[title={Příklad 2}]
Vyřeš, resp. urči: $2\log_3 9-\log_3 3$.
\tcblower
\[\begin{aligned}\log_3 81-\log_3 3 \\ =\log_3 27=3\end{aligned}\]
\end{prikladReseniBox}
\poznamka{Zdůrazni, že $\log_a(u+v)$ nemá žádné obdobné rozkladové pravidlo. Podmínky $u,v>0$ říkej při každém použití.}
\blok{Samostatné procvičování}
\begin{zadaniReseniBox}[title={Samostatná práce 1}]
\begin{enumerate}[label=\alph*)]
\item $\log_{2}8+\log_{2}4$.
\item $\log_{3}9-\log_{3}3$.
\item $\log_{2}16+\log_{2}2$.
\item $\log_{5}25-\log_{5}5$.
\item $\log_{2}32+\log_{2}8$.
\item $\log_{3}27-\log_{3}9$.
\end{enumerate}
\tcblower
\begin{enumerate}[label=\alph*)]
\item $5$.
\item $1$.
\item $5$.
\item $1$.
\item $8$.
\item $1$.
\end{enumerate}
\end{zadaniReseniBox}
\blok{Dobrovolné rozšíření}
\begin{zadaniReseniBox}[title={Rozšíření 1}]
\begin{enumerate}[label=\alph*)]
\item $Vytvoř vlastní úlohu tohoto typu a se spolužákem porovnej podmínky a výsledek.$.
\end{enumerate}
\tcblower
\begin{enumerate}[label=\alph*)]
\item Zkontroluj definice a argumenty..
\end{enumerate}
\end{zadaniReseniBox}
\blok{Úkol na doma}
U každého logaritmu vypiš podmínky argumentu a proveď zkoušku. U rovnic vyznač případně nepřípustné kořeny, u nerovnic odpověz intervalem.
\begin{zadaniReseniBox}[title={Domácí úkol}]
\begin{enumerate}[label=\alph*)]
\item $\log_{2}4+\log_{2}8$.
\item $\log_{3}3-\log_{3}27$.
\item $\log_{5}125+\log_{5}25$.
\item $\log_{2}64-\log_{2}4$.
\item $\log_{3}81+\log_{3}9$.
\item $\log_{2}8-\log_{2}32$.
\item $\log_{5}25+\log_{5}125$.
\item $\log_{2}2-\log_{2}4$.
\end{enumerate}
\tcblower
\begin{enumerate}[label=\alph*)]
\item $5$.
\item $-2$.
\item $5$.
\item $4$.
\item $6$.
\item $-2$.
\item $5$.
\item $-1$.
\end{enumerate}
\end{zadaniReseniBox}
\poznamka{Domácí úkol obsahuje 8 položek; předpokládaná domácí práce přibližně 15--30 minut. Po větších blocích jsou vědomě promíchané dříve probrané typy.}
\end{document}
